\section{Small derivatives of the remote action}
\label{sec:remote-action-jets}
The packet estimate permits a polynomial expense in its logarithm.
Outside the packets we need a stronger conclusion: each local
logarithmic activity, together with its derivatives, must be small.
We obtain the extra power from the local quadratic comparison and
the flat interpolation argument already proved. Throughout this
section the derivative order, block length, patch, forest and tags
are fixed. The path bounds, flat identities and positive precision
estimates retain the hypotheses of the preceding sections.

\subsection{A cubic remainder in a flat chart}
For a local energy address $p$, let $\mathcal E_p(U,V)$ denote its
unmasked energy, with the original predictor tags. Let $I(p)$ be
its set of extracted quadratic rows; $|I(p)|\le2$. Write $f_i(V)$
for the zero-fluctuation residual and $M_i(V)u$ for the transported
cut linear row. Define
\begin{align}
 R_p(U,V,u)&=\mathcal E_p(U,V)
       -\frac12\sum_{i\in I(p)}|f_i(V)+M_i(V)u|^2,
       \label{eq:remote-square-remainder}\\
 F_p(b,u)&=R_p(\exp(u)\operatorname{replica}(\exp b),\exp b,u).
       \label{eq:remote-flat-remainder}
\end{align}
The four address types are the fine curvature, predictor penalty,
coarse curvature with its two mismatch terms, and the null-row
square. For the last two, the parent set is empty. These are the
same unmasked expressions as in Section~\ref{sec:packet-annulus}.

For every fixed $k$, there are a radius $r>0$ independent of the
period and a function $A$ of upper exponent zero such that
\begin{equation}\label{eq:remote-fixed-ball}
 F_p\in\mathcal B_k(A(x),1;(b,u)),\qquad
 \max(\|b\|_\infty,\|u\|_\infty)\le r.
\end{equation}
Here and below the same bound holds at every address. To verify it,
use the four dressed fields and their predictor bounds in
Section~\ref{sec:mean-action-jets}. Each cut path has length
$O(R+1)$, its absolute coefficients have bounded sum, and its
transport derivatives cost a polynomial in $R+1$. The local energy
is a sum of squared curvature rows and inner products of the
coarse stroke with linear guide rows. Its derivatives therefore
cost only another polynomial in $R+1$. Each affine row in
\eqref{eq:remote-square-remainder} is the sum of the residual
already estimated there and a cut linear map applied to $u$.
On the unit ball, the coordinate map $u\mapsto u_e$ has value and
first derivative at most one and all higher derivatives zero.
The product rule bounds its transported image by a polynomial
in $B+1$. Squaring and summing at most two such rows proves
\eqref{eq:remote-fixed-ball}. This is a bound on the unmasked
formula; it makes no small-field assertion on that whole fixed ball.

Put $t=q h^{-1/2}$ with $q>0$ fixed and sufficiently small.
Every coarse field $\exp b$ with $\|b\|_\infty\le3t$ is clean
on all its plaquettes. Indeed the exponential chart gives
$\operatorname{chord}(\exp b_e)\le C\|b_e\|$, and the four-link
triangle inequality gives plaquette chord at most $12Cq h^{-1/2}$.
Choose $q\le(12C)^{-1}$ and use $h^{-1/2}$ no larger than the
coarse flag threshold. This check applies on every period.

On the radius-$3t$ ball, the fields in
\eqref{eq:remote-flat-remainder} have the exact local graft
$U=\exp u\operatorname{replica}V$. The annular comparison then
gives
\begin{equation}\label{eq:remote-probe-value}
 |F_p(b,u)|\le E(x),\qquad E\text{ of upper exponent }-3/2.
\end{equation}
More explicitly, the fine and predictor rows have affine error of
upper exponent $-1$ and amplitude of upper exponent $-1/2$;
the difference of their squares consequently has upper exponent
$-3/2$. Their extracted rows agree with $f_i+M_i u$, since
cleanliness removes the clipping on the complete read stencil.
The coarse curvature remainder and null-row square have upper
exponent $-2$. Each mismatch term is an inner product of a coarse
stroke of upper exponent $-1/2$ with an error of upper exponent
$-1$. Their sum has upper exponent $-3/2$. Taking the sum of
these four nonnegative bounds gives one $E$ for all address types.
These are precisely the estimates proved in the annular argument;
no inequality is differentiated in obtaining
\eqref{eq:remote-probe-value}.

Let $Q\ge0$ have upper exponent $-1/2$, and put $B_0=Q+3t$.
Apply \eqref{eq:action-interpolation} to $F_p$, with $k=j+3$,
value bound $E$ and fixed-ball derivative bound $A$. Its constants
satisfy
\[
 W=1000(E/t^2+AB_0),\qquad
 Y=1000E/t+3WB_0,\qquad
 H=E+YB_0+Yt+Wt^2+At^3.
\]
Their upper exponents are $-1/2,-1,-3/2$, respectively. For example
$E/t^2$ and $AB_0$ have upper exponent $-1/2$, and each of
$YB_0,Yt,Wt^2,At^3$ has upper exponent $-3/2$. Thus
\begin{equation}\label{eq:remote-probe-jets}
 F_p\in\mathcal B_j(H,t^{-1};z)
 \quad(\|z\|\le Q),\qquad H\text{ of upper exponent }-3/2.
\end{equation}
The reciprocal $t^{-1}$ has upper exponent $1/2$ and is paid by
the small derivatives of the chart coordinates in the next step.

\subsection{Returning to a protected remote address}
Suppose the anchor of $p$ is farther than $200B$ from the seeds.
Let $V(y)$ be a coarse variation of derivative scale with upper
exponent $-1/2$, and let $u(y)$ be its coherently compensated fine
increment. Hold the underlying fine configuration $U_*$ fixed,
as in Section~\ref{sec:compensated-variation}. Let $\psi_p(y)$
be the product of the doubled protective cutoffs on the local
read region. If $\psi_p$ vanishes on a neighbourhood of $y_0$,
the protected remainder has zero jets there. Otherwise the closed
support argument bounds the increment on that read region by
$q_0+4a\tau$, with $q_0$ of upper exponent $-1/2$.

After a signed pullback, the relative increment is
$\operatorname{Ad}$ of this increment by the appropriate replica
link. The product rule and the unitary bounds give amplitude
\[
 R_G^2(q_0+4a\tau)\{1+(v+1+v)\}^{j}
\]
for its jets through order $j$; this still has upper exponent
$-1/2$. On tree bonds the increment is identically zero. All other
bonds in the radius-$B$ read region are active because the address
is remote from the seeds. The compensated identity therefore gives,
on a neighbourhood of $y_0$,
\[
 U_* =\exp u(y)\operatorname{replica}V(y)
 \quad\text{on the radius-}\lfloor7B/8\rfloor\text{ patch},
\]
with the signed pullback understood on both sides.

Choose one coarse gauge from $V(y_0)$ on this patch. The coarse
logarithmic coordinates $b(y)$ have jets of upper exponent $-1/2$
by the fixed-gauge construction in Section~\ref{sec:mean-action-jets}.
Conjugate $u(y)$ by this fixed gauge and set it to zero outside
the patch. This preserves its derivative bound. Denote the resulting
pair of arrays by $z(y)$. For one nonnegative $Q$ of upper exponent
$-1/2$, we have $z\in\mathcal B_j(Q,1;y_0)$.

All terms in \eqref{eq:remote-square-remainder} are determined by
this patch. To check the radius, the cut guide and transported
linear rows read within $B/4$ of their address; the local energy
and its predicted coarse field read within $B/2$, including the
fixed predictor reach. Both fit inside $7B/8$ for the previously
chosen geometric threshold. The row transports, predictor tags and
quadratic extraction decisions are unchanged there. Gauge covariance
and the local equality of the fields consequently identify the
actual remainder with $F_p(z(y))$ throughout a neighbourhood.

For a derivative of order $k\le j$, the chain rule has at most
$k!$ partition terms. A term with $m$ blocks is bounded by
$H t^{-m}Q^m$. Hence every derivative of $F_p(z(y))$ is bounded
by
\[
 H\max(1,j!) (1+Qt^{-1})^j.
\]
The second factor has upper exponent zero. We have proved that
the remote unmasked remainder has jets of upper exponent $-3/2$
whenever its protector is not locally zero. The protector has
amplitude one and derivative scale bounded by a function of upper
exponent zero times its support cover. Multiplying it back in
preserves the amplitude exponent $-3/2$, allowing this support
factor in the derivative scale.

\subsection{The remote mean square and other short terms}
The local mean residual $m_p$ has remote value bound of upper
exponent $-1$, proved in Section~\ref{sec:local-mean}. Its jets
have the same exponent. Here is why the scalar estimate suffices
in this application. On a flat probe ball all sectors are clean
and the actual mean agrees with its unmasked expression. The
derivative estimate on the larger \emph{shrinking} ball requires
an additional check, since the intermediate probes need not satisfy
the sector flag conditions. Let $B_0=Q+3t$ as above and write
$A(z)$ for the empty-seed precision at logarithmic probe $z$.
The same-history comparison at the identity supplies
$A(0)\ge r_0 I$ for a fixed $r_0>0$. On $\|z\|\le B_0$,
the transported-row estimates give an entry derivative bound
$C v_1$, where $C$ is fixed and $v_1$ has upper exponent zero.
Integrating this derivative along the segment from $0$ to $z$
gives $|A(z)_{ik}-A(0)_{ik}|\le C v_1B_0$. The matrices have
range at most $B$ and coordinate density at most $\rho$ per site.
Both the row and column sums of this difference are therefore
at most $b_B\rho C v_1B_0$, which tends to zero. The elementary
row-column bound on the Euclidean operator norm follows by
Cauchy--Schwarz:
$\|Hv\|_2^2\le(\sup_i\sum_k|H_{ik}|)
(\sup_k\sum_i|H_{ik}|)\|v\|_2^2$.
Consequently $A(z)\ge(r_0/2)I$ throughout the shrinking ball
for large $x$. Every principal restriction has the same floor.
The exponentially weighted entry estimates remain uniform there;
the inverse and padded-inverse jet bounds proved earlier thus
apply with fixed reduced floor and width. They cost only a
polynomial in $B$ at fixed derivative order.

The empty-seed residual is the zero-fluctuation row residual
minus the transported row applied to its padded-inverse response,
that is $f-MPM^*f$ after the fixed row normalization. All four
factors now have jet bounds of upper exponent zero on this
shrinking ball. The product and row-sum estimates give the required
upper exponent zero bound for the residual there. On the smaller
radius-$3t$ ball every empty-seed probe is clean, so the
same-history stationary comparison applies separately to each
probe and gives the remote value exponent $-1$. No flag condition
was assumed on the entire larger ball. Apply the same interpolation
above with this probe value exponent $-1$.
Then $W,Y,H$ have upper exponents $0,-1/2,-1$. The fixed-gauge
composition again costs only upper exponent zero. Finally the
mean residual depends only on the radius-$24B$ patch and the
unchanged windows; a point farther than $30B$ from the seeds has
the same local formula as the empty-seed comparison. This gives
the claimed exponent $-1$ for every fixed jet order, with no
derivative of a gauge choice or stationary minimizer.

At most two such residuals enter a local mean-square term, so its
jets have upper exponent $-2$. The remote logarithmic action is
$-\theta\beta$ times the sum of this mean-square term and the
protected remainder, where $0\le\theta\le1$ is the fixed channel
weight and $0<\beta\le2h$. Its amplitude therefore has upper
exponent $-1/2$. This uses the same $\beta$ as in the original
action.

The other two short logarithmic terms are the remote normalization
correction and the protected chart logarithm. Their jets have upper
exponent $-1/2$ by Sections~\ref{sec:normalization-jets} and
\ref{sec:haar-chart}. Each term is set to zero on its prescribed
packet neighbourhood; this indicator depends only on the frozen
sector. Consequently the three short families have one common
nonnegative amplitude $b(x)$ of upper exponent $-1/2$, and
derivative scale
\begin{equation}\label{eq:remote-short-price}
 v_0(x)\{1+\operatorname{cov}_B(S_p)\},
 \qquad v_0\text{ of upper exponent zero}.
\end{equation}
Here $S_p$ is the full enlarged read support of the term.
For fixed enlargement $\ell$, these supports lie in radius
$(\ell+1)B$ balls about their anchors and have uniformly bounded
cover number. The number of anchors meeting a given site is at
most $b_{(\ell+1)B}$ times the fixed density of the three families.
Their exponentially cover-weighted site sum is therefore bounded
by this polynomial count times $b(x)$ and a fixed exponential
in the cover bound. It still has upper exponent $-1/2$, uniformly
in the period. This supplies the short-activity part of the
logarithmic price estimate.
